\subsection{Le modèle du microscope en optique géométrique}

Le microscope est un instrument d'optique conçu pour l'observation d'objets 
de dimensions micrométriques situés à \textbf{distance finie}. 
Son objectif principal est de fournir une \textbf{image virtuelle} très agrandie de l'objet, 
permettant ainsi d'atteindre des détails inaccessibles à l'œil nu.

Dans un modèle réduit, 
le microscope est assimilé à un système de deux lentilles minces convergentes de même axe optique :

\begin{itemize}
    \item \textbf{L'objectif ($L_1$)} : C'est la lentille située du côté de l'objet. 
    Elle possède une distance focale $f'_1$ très courte (de l'ordre de quelques millimètres). 
    Son rôle est de projeter de l'objet $AB$ une image intermédiaire $A_1B_1$ réelle et renversée. 
    Pour cela, l'objet doit être placé légèrement au-delà du foyer objet $F_1$.
    
    \item \textbf{L'oculaire ($L_2$)} : Située du côté de l'œil, 
    cette lentille de distance focale $f'_2$ plus importante joue le rôle d'une loupe. 
    Elle permet d'examiner l'image intermédiaire $A_1B_1$ en en donnant une image finale virtuelle, 
    redressée par rapport à $A_1B_1$ 
    (et donc renversée par rapport à l'objet initial $AB$).
\end{itemize}

Une caractéristique fondamentale du microscope est son \textbf{intervalle optique}, noté $\Delta$. 
Il correspond à la distance algébrique séparant le foyer image de l'objectif $F'_1$ du foyer objet de l'oculaire $F_2$ :
\begin{equation}
    \Delta = \overline{F'_1 F_2}
\end{equation}

% Pense à insérer ici ton schéma avec l'environnement figure
% \begin{figure}[h]
%    \centering
%    \includegraphics[width=0.8\textwidth]{ton_schema.png}
%    \caption{Schéma optique du microscope réduit.}
%    \label{fig:schema_optique}
% \end{figure}

\subsection{Condition de vision sans fatigue : mise au point à l'infini}

Pour un œil normal (emmétrope), l'observation d'un objet n'engendre aucune fatigue nerveuse ou musculaire lorsque celui-ci est situé à l'infini (point appelé \textit{Punctum Remotum}). Dans cette configuration, les muscles ciliaires sont au repos et l'œil n'accommode pas.

Pour réaliser cette condition avec un microscope, l'image finale $A'B'$ donnée par l'instrument doit être rejetée à l'infini. Compte tenu du modèle à deux lentilles :
\begin{itemize}
    \item L'objectif $L_1$ donne de l'objet $AB$ une image intermédiaire $A_1B_1$.
    \item L'oculaire $L_2$ agit comme une loupe pour observer $A_1B_1$.
\end{itemize}

Par conséquent, pour que l'image finale soit à l'infini, l'image intermédiaire $A_1B_1$ doit impérativement se situer dans le \textbf{plan focal objet de l'oculaire} ($F_2$).

La manipulation de mise au point consiste à faire varier la distance entre l'objet et l'objectif 
(à l'aide de la crémaillère ou de la vis micrométrique) 
afin de faire coïncider l'image $A_1B_1$ avec ce plan $F_2$. 
Bien que l'œil humain soit capable d'accommoder pour voir net une image située entre son 
\textit{Punctum Remotum} (infini) et son 
\textit{Punctum Proximum} (considéré à $25 cm$ en optique géométrique), 
la visée à l'infini est privilégiée pour une observation prolongée.